\documentclass[10pt,a4paper]{amsart}
\usepackage[margin=19mm]{geometry}
\usepackage{amsmath,amssymb,mathtools}
\usepackage[scr=rsfs,cal=euler]{mathalfa}
\usepackage{xfrac}
\usepackage[unicode,hidelinks,bookmarks=false]{hyperref}
\newcommand{\ie}{i.e.\ }
\newcommand{\cf}{cf.\ }
\newcommand{\resp}{respectively\ }
\newcommand{\Subsection}{\subsection}
\providecommand{\nobreakdash}{\nobreak\hbox{-}}
\setlength{\emergencystretch}{2em}
\multlinegap0pt
\usepackage{amssymb}
\usepackage{enumerate}
\usepackage{xcolor}
\usepackage{subfig}
\usepackage{mathtools}
\newtheorem{lemma}{Lemma}
\newcommand{\mchi}{\raisebox{0pt}[1ex][1ex]{$\chi$}}
\newcommand{\veph}{\hat\varepsilon }
\title{Nonradial Quenching Profile for a MEMS Model}
\author{Hsuan-Lin Liao, Van Tien Nguyen}
\date{Source: arXiv:2510.15246v2 (CC BY 4.0)}
\begin{document}
\maketitle

\noindent\emph{Source and licence.} Nonradial Quenching Profile for a MEMS Model. Version arXiv:2510.15246v2 (CC BY 4.0), distributed by arXiv under the Creative Commons Attribution 4.0 International licence, http://creativecommons.org/licenses/by/4.0/, as recorded in the arXiv licence metadata for this exact version and shown on the abstract page https://arxiv.org/abs/2510.15246v2. Full licence notice: \texttt{licenses/arxiv-2510.15246-CC-BY-4.0.txt}.\par\smallskip

\noindent\textit{Editorial scope.} Counted unit: the article's lemma Lemma4.1 (source0.tex lines 744--782). The article's complete original proof of it is delivered with it (source0.tex lines 783--827, one \texttt{proof} environment of 45 lines). No mathematics is written, repaired or re-derived: the statement and the proof are the article's own lines, in the article's own notation, and no retained line is substituted or re-flowed.  Retained uncounted as the source-local context the proof needs: lines 537--539; lines 562--568; lines 570--573.  Nothing was omitted from any retained span, so the package carries no omission record.  No retained line contains a citation, so the wrapper carries no bibliography and no bibliography entry.  No retained line contains a figure environment, an image-inclusion command or drawing code, so no image is delivered and the package declares no asset dependency.  Editorial formatting only, in addition: an AMS \texttt{amsart} preamble that restates the article's own declarations and macros used by the retained lines, namely the environments \texttt{lemma} on separate counters, together with the standard packages those lines need.  The retained environments print in this document as Lemma 1 in order of appearance, whereas the article prints the same items with the numbering of its own full document.  The complete native source of the article as distributed in the arXiv e-print for this exact version is bundled unchanged as \texttt{sources/187241/source0.tex}.
\begin{align}
    \veph(y,s)=\sum_{i+j \leq 8}\veph_{ij}(s)H_{ij}(y)+\veph_-(y,s),\label{decomoin} \quad \quad \veph_- \perp_{L^2_\rho} \{H_{ij}\}_{i+j \leq 8},
\end{align}

\begin{align}
    &\left( \sum_{i+j \leq 7} |\veph_{ij}(s)|^2 \right)^{\frac{1}{2}} \leq A_1 e^{-\frac{32}{12}s},\label{in1}\\
    &|\veph_{ij}(s)| \leq A_1e^{-\frac{32}{12}s}, \text{ for } i+j =8, \label{in2}\\
    &\|\veph_-(s)\|_{L^2_{\rho}(\mathbb{R}^2)} \leq A_1 e^{-\frac{32}{12}s} \label{in3},\\
    &\|\partial_i\veph_-(s)\|_{L^2_{\rho}(\mathbb{R}^2)} \leq A_2 e^{-\frac{32}{12}s}, \quad i=1,2,\label{in4}\\
    &\|\partial^2_{ij}\veph_-(s)\|_{L^2_{\rho}(\mathbb{R}^2)} \leq A_3 e^{-\frac{32}{12}s}, \quad i,j=1,2, \label{in5}
\end{align}

\begin{align}
    &\|(y \cdot \nabla)^k\varepsilon\|_{\flat} \leq B_{k+1} e^{-\frac{32}{12}s}, \quad k=0,1,2,\label{int1}\\
   \quad  &\|(z \cdot \nabla)^k \eta \|_{\natural} \leq C_{k+1} e^{-\frac{5}{12}s}, \label{bootstrap: int 2} \quad k=0,1,2,
\end{align}

\begin{lemma}
\label{Lemma4.1}
    Let \(s_0 = s_0(A_1, A_2,A_3,K,B_1,B_2,B_3,C_1,C_2,C_3,D_1) \gg 1\), let \(\varepsilon(s) \in S(s)\) be a solution for \(s \in [s_0, s_1]\), with fixed \(s_1 > s_0\). Then we have the following estimates:\\
(i) \textup{(Local \(L^{\infty}\)-estimate for \(\veph_-\).) }
    \begin{equation*}
        \|\veph_-(s)\|_{L^{\infty}(|y| \leq 100K)} \lesssim C(K) \, e^{-\frac{32}{12}s},
    \end{equation*}
    where \(C(K)>0\) is a constant depending only on \(K\) and \(\veph_-\) is defined in \ref{decomoin}.\\
(ii) \textup{(Pointwise estimates for \(\varepsilon\) and \(\eta\)}.) 
We have the pointwise estimates
\begin{align*}
    |\varepsilon(y,s)| \lesssim e^{-\frac{32}{12}s}|y|^9, \text{ for } |y| \geq 2K, \quad |\eta(z,s)| \lesssim e^{-\frac{5}{12}s}|z|^{4.005}, \text{ for } |z| \geq 2K.
\end{align*}
(iii) \textup{(Nonlinear estimate for \(\veph\) in \(L^2_{\rho}(\mathbb{R}^2)\).)} 
The nonlinear term corresponding to \(\veph\) is defined as
\begin{align*}
        NL(\veph)=\sum_{k=2}^{\infty}\binom{-2}{k}\frac{\veph^k}{\mathcal{P}_{\theta}^{k+2}},
    \end{align*}
    where 
    \begin{align*} \binom{-2}{k}=\frac{(-2) \cdot (-3) \cdots(-2-k+1)}{k!}, \quad k \geq 2,
    \end{align*}
    and satisfies the estimate
    \begin{align*}
        NL(\veph)=o(se^{-3s}), \quad \text{ in } L^2_{\rho}(\mathbb{R}^2).
\end{align*}
(iv) \textup{(Exterior nonlinear estimate)}
The exterior nonlinear term is defined as
\begin{align*}
    NL^{ex}(\eta)=NL(\eta)\left(1-\mchi_{K}(ze^{-\frac{5}{49}s})\right),
\end{align*}
where
\begin{align*}
    NL(\eta):=\left(\frac{1}{\mathcal{P}_{\theta}^2}-\frac{1}{(\mathcal{P}_{\theta}+\eta)^2}\right),
\end{align*}
and satisfies the estimate
\begin{align*}
    |NL^{ex}(\eta)| \lesssim e^{-\frac{11}{147}s}
\end{align*}
\end{lemma}

\begin{proof}(i) Note that by Sobolev inequality, we have
    \begin{align*}
      \|\veph_-(s)\|^2_{L^{\infty}(|y| \leq 100K)} \leq C\left(\|D^2 \veph_-(s)\|^2_{L^2(|y| \leq 50K)}+\|D \veph_-(s)\|^2_{L^2(|y| \leq 50K)}+\|\veph_-(s)\|^2_{L^2(|y| \leq 50K)}\right),
    \end{align*}
    where \(C>0\) is the Sobolev constant. By the bootstrap condition (\ref{in3}), we have:
   \begin{align*}
    \|\veph_-(s)\|_{L^2(|y| \leq 50K)}& \lesssim C(K) \|\veph_-\|_{L^2_{\rho}(\mathbb{R}^2)} \lesssim C(K)\,e^{-\frac{32}{12}s}.
   \end{align*}
Next, by the bootstrap condition \ref{in4} and \ref{in5}, we have
   \begin{align*}
      &\|D \veph_-(s)\|_{L^2(|y| \leq 50K)} \lesssim C(K)\,e^{-\frac{32}{12}s}, \quad \|D^2 \veph_-(s)\|_{L^2(|y| \leq 50K)} \lesssim C(K)\,e^{-\frac{32}{12}s},
   \end{align*}
hence, we have
   \begin{align*}
      \|\veph_-(s)\|_{L^{\infty}(|y| \leq 100 K)} \lesssim  C(K)\,e^{-\frac{32}{12}s}.
   \end{align*} 
This concludes the local \(L^{\infty}\) estimate of \(\veph_-\).\\
 (ii) This follows directly from bootstrap conditions \ref{int1} and \ref{bootstrap: int 2} combined with Sobolev inequality.\\
(iii)
Note that by (ii), we have
\begin{align}
\label{inq:ptwise1}
    |\veph(y,s)| \leq M e^{-\frac{32}{12}s}|y|^9, \quad |y| \geq 2K,
\end{align}
moreover, by the local \(L^{\infty}\) estimate for \(\veph_-\), together with the bootstrap condition \ref{in1} and \ref{in2}, we have
\begin{align}
\label{ptwise local veph 0}
    \|\veph(y,s)\|_{L^{\infty}(|y|\leq 50K)} \leq Me^{-\frac{32}{12}s}.
\end{align}
for some constant \(M=M(A_1,A_2,A_3,K,B_1,B_2,B_3)>0\). Thus, we have, for \(k \geq 2\),
    \begin{align*}
        \|-(\veph^k/\mathcal{P}_{\theta}^{k+2})\|_{L^2_{\rho}(\mathbb{R}^2)}& \leq \|\veph^k\|_{L^2_{\rho}(\mathbb{R}^2)}\\
        &=\left(\int_{|y| \leq 2K}|\veph|^{2k}\rho(y)\,dy+\int_{|y| \geq 2K}|\veph|^{2k}\rho(y)\,dy\right)^{1/2} \leq 2M^k\,e^{-\frac{32}{12}ks},
    \end{align*}
    and hence
    \begin{align*}
        \|NL(\veph)\|_{L^2_{\rho}(\mathbb{R}^2)} &\leq \sum_{k=2}^{\infty}(k+1) \|\veph^k\|_{L^2_{\rho}(\mathbb{R}^2)} \leq\sum_{k=2}^{\infty}(k+1)(2M^ke^{-\frac{32}{12}ks})\ll se^{-3s}.
    \end{align*}
    This proves the nonlinear estimate for \(\veph\) in \(L^2_{\rho}(\mathbb{R}^2)\).\\
   (iv) Note that 
    \begin{align*}
        |NL^{ex}(\eta)| \lesssim\left( \frac{1}{\mathcal{P}_{\theta}^2}+\frac{1}{(\mathcal{P}_{\theta}-|\eta|)^2}\right)\left(1-\mchi_{K}(ze^{-\frac{5}{49}s})\right)\lesssim \frac{1}{\mathcal{P}_{\theta}^2}\left(1-\mchi_{K}(ze^{-\frac{5}{49}s})\right) \lesssim e^{-\frac{11}{147}s},
    \end{align*}
    where we used the fact that \(|\eta^{ex}(z,s)/\mathcal{P}_{\theta}(z,s)| \ll 1\). This concludes the exterior nonlinear estimate.
\end{proof}

\end{document}
